\chapter{Esercizi: Confronti, cicli e membri di classe in Java}
\label{chap:esercizi-confronti-cicli-java}

\begin{exercise}
\label{ex:fi-esercizi-java-cicli-stringhe-equals}
Perché l'uguaglianza del contenuto di due stringhe deve essere verificata
con \texttt{equals()} anziché con \texttt{==}?
\end{exercise}

L'operatore \texttt{==} confronta l'identità dei riferimenti, mentre
\texttt{String.equals()} confronta i caratteri delle stringhe.
Due stringhe possono quindi avere lo stesso testo e appartenere a
oggetti differenti
(\cref{sec:fi-java-cicli-stringhe-uguaglianza}).

\begin{exercise}
\label{ex:fi-esercizi-java-cicli-ignorecase}
Qual è la differenza tra \texttt{equals()} e
\texttt{equalsIgnoreCase()}?
\end{exercise}

\texttt{equals()} richiede che il contenuto delle stringhe coincida;
\texttt{equalsIgnoreCase()} trascura le differenze tra maiuscole e
minuscole secondo le regole di confronto di \texttt{String}.
Non rimuove in generale gli accenti e non implementa un ordinamento
linguistico
(\cref{sec:fi-java-cicli-stringhe-uguaglianza}).

\begin{exercise}
\label{ex:fi-esercizi-java-cicli-stringhe-uguale-riferimento}
Che cosa verifica \texttt{s1 == s2} quando entrambe le variabili
sono riferimenti a oggetti \texttt{String}?
\end{exercise}

Verifica che le variabili contengano lo stesso riferimento, cioè
che individuino il medesimo oggetto, oppure che entrambe valgano
\texttt{null}. Non confronta il contenuto dei due testi
(\cref{sec:fi-java-cicli-stringhe-uguaglianza}).

\begin{exercise}
\label{ex:fi-esercizi-java-cicli-stringhe-identita-valore}
Come può risultare vera \texttt{s1.equals(s3)} mentre
\texttt{s1 == s3} risulta falsa?
\end{exercise}

Le due variabili possono riferirsi a oggetti \texttt{String}
distinti che contengono esattamente gli stessi caratteri.
Il primo confronto verifica il valore, il secondo l'identità
(\cref{fig:fi-java-cicli-stringhe-riferimenti}).

\begin{exercise}
\label{ex:fi-esercizi-java-cicli-compareto}
Quale risultato restituisce \texttt{compareTo()} e come se ne
interpreta il segno?
\end{exercise}

Restituisce un intero negativo se il ricevente precede l'argomento
nell'ordinamento, zero se i testi coincidono e positivo se lo segue.
È necessario controllare il segno, non aspettarsi i valori esatti
\(-1\), \(0\) o \(1\)
(\cref{sec:fi-java-cicli-stringhe-ordine}).

\begin{exercise}
\label{ex:fi-esercizi-java-cicli-lessicografico}
In che cosa l'ordinamento lessicografico di \texttt{String}
si distingue dall'ordinamento alfabetico usuale?
\end{exercise}

Il confronto segue l'ordine delle unità di codice UTF-16, non le
regole di un dizionario italiano. Distingue maiuscole e minuscole
e considera anche cifre, spazi e simboli; in caso di prefisso
comune, la stringa più corta precede quella più lunga
(\cref{sec:fi-java-cicli-stringhe-ordine}).

\begin{exercise}
\label{ex:fi-esercizi-java-cicli-oggetti-operatore}
Che cosa verifica \texttt{==} quando è applicato a due variabili
riferimento?
\end{exercise}

Verifica l'identità degli oggetti indicati, oppure che entrambi i
riferimenti siano nulli. Non stabilisce in generale se gli oggetti
abbiano lo stesso valore logico
(\cref{def:fi-java-cicli-uguaglianza-oggetti}).

\begin{exercise}
\label{ex:fi-esercizi-java-cicli-equals-object}
Perché \texttt{equals()} è disponibile per gli oggetti, ma
viene spesso ridefinito nelle singole classi?
\end{exercise}

È dichiarato nella classe \texttt{Object}, la cui implementazione
verifica l'identità. Una classe può ridefinirlo per stabilire
un criterio di equivalenza coerente con il significato dei propri
oggetti, come avviene per le stringhe
(\cref{sec:fi-java-cicli-oggetti}).

\begin{exercise}
\label{ex:fi-esercizi-java-cicli-comparable}
Perché non tutti gli oggetti possono essere confrontati con
\texttt{compareTo()}?
\end{exercise}

A differenza di \texttt{equals()}, \texttt{compareTo()} non è
dichiarato nella classe \texttt{Object}: occorre una classe che
definisca un ordinamento naturale, per esempio implementando
\texttt{Comparable}. Il criterio di ordinamento dipende dal tipo
(\cref{sec:fi-java-cicli-oggetti}).

\begin{exercise}
\label{ex:fi-esercizi-java-cicli-catena-if}
Quando è appropriata una catena
\texttt{if}/\texttt{else if}/\texttt{else}?
\end{exercise}

Quando si deve scegliere un solo comportamento fra più alternative,
anche in base a condizioni diverse o a intervalli di valori.
Le condizioni vengono verificate nell'ordine e, alla prima vera,
i rami successivi sono esclusi
(\cref{sec:fi-java-cicli-if-catene}).

\begin{exercise}
\label{ex:fi-esercizi-java-cicli-soglie}
Perché, in una classificazione mediante soglie del tipo
\texttt{x >= soglia}, è importante ordinare correttamente i controlli?
\end{exercise}

Con una catena di alternative il primo ramo vero viene eseguito.
Se la soglia più bassa compare per prima, può intercettare anche
valori che dovrebbero rientrare in fasce superiori. Nel caso
presentato, le soglie vanno disposte in ordine decrescente
(\cref{sec:fi-java-cicli-if-catene}).

\begin{exercise}
\label{ex:fi-esercizi-java-cicli-else-sospeso}
Che cos'è il problema dell'\emph{else sospeso} e a quale
\texttt{if} viene associata la clausola \texttt{else}?
\end{exercise}

Quando più \texttt{if} sono annidati senza graffe, l'indentazione
può suggerire un'associazione diversa da quella reale. Java
associa \texttt{else} all'\texttt{if} più vicino ancora privo
di \texttt{else}; usare blocchi espliciti elimina l'ambiguità
(\cref{sec:fi-java-cicli-if-catene}).

\begin{exercise}
\label{ex:fi-esercizi-java-cicli-switch}
A che cosa serve \texttt{switch} e quando può risultare più
leggibile di una catena di \texttt{if}/\texttt{else if}?
\end{exercise}

Seleziona un ramo in base al valore di un'espressione, fra un
insieme di casi discreti. È spesso più chiaro quando si confronta
un selettore con diversi valori specifici; non è adatto a
sostituire direttamente controlli su intervalli arbitrari
(\cref{sec:fi-java-cicli-switch}).

\begin{exercise}
\label{ex:fi-esercizi-java-cicli-fallthrough}
Che cosa può accadere in uno \texttt{switch} tradizionale
se manca \texttt{break} alla fine di un \texttt{case}?
\end{exercise}

L'esecuzione può proseguire negli enunciati del caso seguente
senza verificare un'altra etichetta: si tratta del
\emph{fall-through}. Il comportamento è utile per raggruppare
più casi, ma è un errore quando i rami devono rimanere esclusivi
(\cref{sec:fi-java-cicli-switch}).

\begin{exercise}
\label{ex:fi-esercizi-java-cicli-switch-expression}
Quali differenze introduce lo \texttt{switch} usato come
espressione a partire da Java 14?
\end{exercise}

Uno \texttt{switch} può produrre un valore da assegnare o
restituire. La sintassi con \texttt{->} non produce
\emph{fall-through} e la compilazione richiede che siano coperti
tutti i possibili valori del selettore; non occorrono i
\texttt{break} per distinguere i rami
(\cref{sec:fi-java-cicli-switch-expression}).

\begin{exercise}
\label{ex:fi-esercizi-java-cicli-boolean}
Che cosa rappresenta il tipo primitivo \texttt{boolean}?
\end{exercise}

È il tipo dei valori logici \texttt{true} e \texttt{false}.
Le espressioni relazionali e i metodi predicativi producono
valori di tale tipo; in Java un intero non vale automaticamente
come condizione logica
(\cref{def:fi-java-cicli-boolean}).

\begin{exercise}
\label{ex:fi-esercizi-java-cicli-flag}
Che cos'è una variabile booleana utilizzata come \emph{flag}?
\end{exercise}

È una variabile che memorizza una condizione o uno stato con due
possibilità, per esempio se la lettura dell'input debba
continuare. Il suo valore può governare direttamente un ciclo
(\cref{sec:fi-java-cicli-boolean,sec:fi-java-cicli-while-annidati}).

\begin{exercise}
\label{ex:fi-esercizi-java-cicli-predicativi}
Che cosa sono i metodi predicativi e perché spesso iniziano
con \texttt{is} oppure \texttt{has}?
\end{exercise}

Sono metodi il cui risultato è di tipo \texttt{boolean} e che
verificano una proprietà. I prefissi \texttt{is} e \texttt{has}
esprimono la natura della domanda. Ne sono esempi
\texttt{isDigit()} di \texttt{Character} e
\texttt{hasNextInt()} di \texttt{Scanner}
(\cref{sec:fi-java-cicli-boolean}).

\begin{exercise}
\label{ex:fi-esercizi-java-cicli-cortocircuito}
Che cosa significa \emph{valutazione a cortocircuito} per
\texttt{\&\&} e \texttt{||}?
\end{exercise}

Java valuta il secondo operando soltanto se il primo non basta
a determinare il risultato. Con \texttt{\&\&} non valuta il
secondo se il primo è falso; con \texttt{||} non lo valuta se
il primo è vero. Ciò può evitare anche chiamate su riferimenti
nulli (\cref{sec:fi-java-cicli-logici}).

\begin{exercise}
\label{ex:fi-esercizi-java-cicli-relazionale-errore}
Perché \texttt{0 <= x <= 1} non è valido in Java e come si
esprime correttamente la condizione?
\end{exercise}

La prima relazione \texttt{0 <= x} produce un \texttt{boolean},
che non può essere confrontato numericamente con \texttt{1}.
Si usa \texttt{0 <= x \&\& x <= 1}, combinando due
espressioni relazionali booleane
(\cref{sec:fi-java-cicli-precedenze}).

\begin{exercise}
\label{ex:fi-esercizi-java-cicli-enum}
Che cos'è un \texttt{enum} e perché è preferibile a codici
interi o stringhe quando l'insieme dei valori è fissato?
\end{exercise}

È un tipo speciale di classe che definisce un insieme chiuso di
costanti. Il compilatore impedisce di assegnare valori di tipi
incompatibili, mentre nomi significativi e supporto per
\texttt{switch} migliorano la leggibilità
(\cref{sec:fi-java-cicli-enum}).

\begin{exercise}
\label{ex:fi-esercizi-java-cicli-enum-limiti}
Quali limiti presenta un tipo \texttt{enum}?
\end{exercise}

L'insieme delle costanti è fissato nella definizione e non può
essere ampliato dinamicamente durante l'esecuzione. Un valore
esterno, come una stringa di input, deve essere interpretato
esplicitamente; essendo un tipo riferimento, una variabile
enumerativa può inoltre contenere \texttt{null}
(\cref{sec:fi-java-cicli-enum}).

\begin{exercise}
\label{ex:fi-esercizi-java-cicli-switch-enum}
Come si utilizzano i valori di un \texttt{enum} nei casi
\texttt{case} di uno \texttt{switch}?
\end{exercise}

Se il selettore è di un tipo enumerativo, le etichette della
forma tradizionale utilizzano il nome semplice della costante:
per \texttt{Stagione}, il caso si indica con
\texttt{case ESTATE:}, senza il prefisso
\texttt{Stagione.} (\cref{sec:fi-java-cicli-enum}).

\begin{exercise}
\label{ex:fi-esercizi-java-cicli-while-do}
Qual è la differenza essenziale fra \texttt{while} e
\texttt{do}/\texttt{while}?
\end{exercise}

Il ciclo \texttt{while} verifica la condizione prima del corpo,
che può quindi non essere mai eseguito. Nel
\texttt{do}/\texttt{while} il controllo avviene dopo il corpo,
che viene eseguito almeno una volta
(\cref{sec:fi-java-cicli-while,sec:fi-java-cicli-do-while}).

\begin{exercise}
\label{ex:fi-esercizi-java-cicli-investimento-terminazione}
Quali ipotesi garantiscono la terminazione del ciclo che cerca
il raddoppio di un investimento?
\end{exercise}

Nel modello matematico con capitalizzazione annuale, un saldo
iniziale positivo e un tasso costante strettamente positivo fanno
crescere il saldo fino a superare qualsiasi soglia finita. Con tasso
nullo e obiettivo maggiore del saldo iniziale, il ciclo non
terminerebbe. In un programma in virgola mobile, per tassi estremamente piccoli,
l'arrotondamento può tuttavia impedire l'aumento effettivo del saldo
(\cref{sec:fi-java-cicli-investimento,sec:fi-java-precisione}).

\begin{exercise}
\label{ex:fi-esercizi-java-cicli-for-tre-parti}
Quali sono le tre parti dell'intestazione di un ciclo
\texttt{for} e quale ruolo hanno?
\end{exercise}

L'inizializzazione prepara il ciclo ed è eseguita una sola volta;
la condizione viene controllata prima di ogni iterazione;
l'aggiornamento viene eseguito dopo ciascun corpo, prima del
controllo successivo
(\cref{sec:fi-java-cicli-for}).

\begin{exercise}
\label{ex:fi-esercizi-java-cicli-for-annidati}
Come funzionano due cicli \texttt{for} annidati e quando
sono utili?
\end{exercise}

Per ogni iterazione del ciclo esterno viene eseguito il ciclo
interno secondo i suoi limiti. Sono adatti, per esempio, a
elaborare righe e colonne di una tabella; il costo complessivo
dipende dal numero di iterazioni di entrambi
(\cref{sec:fi-java-cicli-for-annidati,chap:complessita}).

\begin{exercise}
\label{ex:fi-esercizi-java-cicli-sentinella}
Che cos'è un valore \emph{sentinella} nella lettura iterativa
dell'input?
\end{exercise}

È un valore riservato per segnalare la conclusione della sequenza,
come \texttt{0} fra numeri o \texttt{"Q"} fra stringhe.
Non viene trattato come un elemento ordinario dei dati da
elaborare
(\cref{sec:fi-java-cicli-while-annidati}).

\begin{exercise}
\label{ex:fi-esercizi-java-cicli-break-continue}
Che cosa fanno le istruzioni \texttt{break} e
\texttt{continue} all'interno dei cicli?
\end{exercise}

\texttt{break}, senza etichetta, termina il ciclo più interno;
\texttt{continue}, senza etichetta, passa all'iterazione successiva.
Nel \texttt{for} la prosecuzione include prima la fase di
aggiornamento. Un uso eccessivo può ridurre la chiarezza
del flusso di controllo
(\cref{sec:fi-java-cicli-while-annidati,sec:fi-java-cicli-for}).

\begin{exercise}
\label{ex:fi-esercizi-java-cicli-scopiazione}
Che cosa significa che una variabile è visibile solo entro
un determinato blocco?
\end{exercise}

Il suo nome può essere usato solo nella regione di codice
individuata dalle regole di ambito. Quando si esce da un blocco
che contiene la dichiarazione, la variabile non può più essere
richiamata dal codice successivo
(\cref{sec:fi-java-cicli-scope}).

\begin{exercise}
\label{ex:fi-esercizi-java-cicli-inizializzazione-locali}
Perché leggere una variabile locale non inizializzata produce
un errore di compilazione in Java?
\end{exercise}

Prima della lettura il compilatore deve poter stabilire che
la variabile ha ricevuto sicuramente un valore lungo ogni
percorso di esecuzione pertinente. Diversamente dai campi,
le variabili locali non ricevono automaticamente un valore
iniziale predefinito
(\cref{sec:fi-java-cicli-scope}).

\begin{exercise}
\label{ex:fi-esercizi-java-cicli-static-differenza}
Qual è la differenza tra un metodo statico e un metodo di istanza?
\end{exercise}

Un metodo statico appartiene alla classe, non ha il ricevente
implicito \texttt{this} e può essere invocato tramite il nome
della classe. Un metodo di istanza opera su un oggetto specifico
e può accedere direttamente ai suoi campi
(\cref{sec:fi-java-cicli-metodi-statici}).

\begin{exercise}
\label{ex:fi-esercizi-java-cicli-static-campi}
Perché un contatore progressivo comune a tutti gli oggetti
deve essere dichiarato \texttt{static}?
\end{exercise}

Un campo di istanza ha una copia distinta per ogni oggetto:
inizializzandolo a zero nel costruttore, ogni oggetto ripartirebbe
dallo stesso numero. Il campo statico possiede invece una copia
condivisa, aggiornata da ogni nuova costruzione
(\cref{sec:fi-java-cicli-campi-statici}).

\begin{exercise}
\label{ex:fi-esercizi-java-cicli-scope-lifetime}
Che differenza c'è tra ambito di visibilità e ciclo di vita
di una variabile?
\end{exercise}

L'ambito determina dove il nome è utilizzabile nel codice;
il ciclo di vita riguarda il tempo per cui il relativo dato
esiste durante l'esecuzione. Per esempio, un campo statico
può permanere per tutta la vita della classe anche quando
non viene utilizzato da un determinato metodo
(\cref{sec:fi-java-cicli-scope,tab:fi-java-cicli-variabili}).

\begin{exercise}
\label{ex:fi-esercizi-java-cicli-gc}
Che cos'è il \emph{garbage collector} e perché normalmente
non si deallocano manualmente gli oggetti in Java?
\end{exercise}

La JVM gestisce automaticamente la memoria degli oggetti e può
recuperare quella occupata dagli oggetti non più raggiungibili.
Il programmatore non deve quindi liberarli esplicitamente;
la raccolta non ha però un momento garantito
(\cref{sec:fi-java-cicli-memoria}).

\begin{exercise}
\label{ex:fi-esercizi-java-cicli-gc-tipi}
Quali sono le principali caratteristiche di Serial, Parallel,
G1, ZGC e Shenandoah?
\end{exercise}

Serial GC privilegia la semplicità; Parallel GC è orientato al
throughput; G1 cerca un equilibrio fra volume di lavoro e
pause; ZGC e Shenandoah sono progettati per ridurre la latenza
della raccolta. Disponibilità e prestazioni dipendono dalla JVM
e dalla versione utilizzata
(\cref{sec:fi-java-cicli-memoria}).

\begin{exercise}
\label{ex:fi-esercizi-java-cicli-gc-default}
Qual è il garbage collector predefinito nelle comuni
configurazioni HotSpot a partire da Java 9, e come può
essere richiesto esplicitamente?
\end{exercise}

Il predefinito è in genere G1, che può essere selezionato con
l'opzione \texttt{-XX:+UseG1GC}. Non si tratta di una regola
universale per tutte le implementazioni e configurazioni della
piattaforma Java
(\cref{sec:fi-java-cicli-memoria}).

\begin{exercise}
\label{ex:fi-esercizi-java-cicli-gc-criteri}
Da quali fattori dipende la scelta del garbage collector
per un'applicazione?
\end{exercise}

Dipende dalle dimensioni e dal comportamento dell'heap, dalle
risorse disponibili e dalle esigenze di throughput e latenza.
La scelta richiede misurazioni nell'ambiente reale; non esiste
un raccoglitore migliore in ogni situazione
(\cref{sec:fi-java-cicli-memoria}).

\begin{exercise}
\label{ex:fi-esercizi-java-cicli-shadowing}
Che cosa significa \emph{shadowing} di una variabile e come
si accede a un campo nascosto da un parametro omonimo?
\end{exercise}

Un nome locale o un parametro possono nascondere un campo
omonimo della classe senza generare un errore. In un metodo
di istanza si usa \texttt{this.nome} per accedere al campo;
per un campo statico si può usare \texttt{NomeClasse.nome}
(\cref{def:fi-java-cicli-shadowing}).

\begin{exercise}
\label{ex:fi-esercizi-java-cicli-private-public}
Come cambia l'accessibilità di un membro \texttt{private}
rispetto a un membro \texttt{public}?
\end{exercise}

Un membro \texttt{private} è accessibile dal codice della
classe che lo dichiara, anche su altre sue istanze. Un membro
\texttt{public} è accessibile anche da codice esterno nei
contesti consentiti, usando un riferimento a un oggetto o
il nome della classe, secondo che sia di istanza o statico
(\cref{sec:fi-java-cicli-shadowing}).
